\section{Real moments and schedules}\label{sec:real-schedules}

We establish both directions between compact-orbit addition and
schedules.  From a real moment solution we construct a schedule whose
switch roots form a basis.  Conversely, a schedule gives a real moment
solution; this implication will be used for the block decompositions
in the next section.  Neither statement assumes a root-lattice
condition.  The construction works in type $D$ and in coordinate
permutation systems, with orthogonal central coordinates carried along
by ordinary addition.

\subsection{Conventions and the path problem}

All roots have squared length two.  In type $D_l$ we use
\[
 \Phi=\{\pm e_i\pm e_j:i\ne j\},\qquad
 \mathcal C=\{x_1\ge\cdots\ge x_{l-1}\ge |x_l|\}.
\]
The root lattice consists of the integral vectors of even coordinate
sum, and the weight lattice consists of the vectors whose coordinates
are either all integral or all half-odd-integral.  We include
$D_2=A_1\times A_1$.  In a permutation system the roots are $e_i-e_j$;
the common-coordinate direction is central.  More generally, whenever
a root subsystem has an orthogonal central complement, we apply the
semisimple arguments after subtracting that complement and retain its
ordinary vector-addition equation.  Its coordinates need not be
integrally normalized.

We use addition labels $\lambda,\mu,\xi$, where $\xi=\nu^*$ is the
dominant representative of $-\nu$.  Thus the root-lattice condition in
Theorem~\ref{thm:main} is equivalent to
$\lambda+\mu-\xi\in Q$.  Let $W$ be the Weyl group.  A \emph{shape}
ending at $\mu$ is a continuous polygonal path
$\eta:[0,1]\longrightarrow\R^l$ with $\eta(0)=0$, $\eta(1)=\mu$,
and every increment in $\mathcal C$.  It is \emph{regular} if all its
slopes lie in the open chamber.

\begin{definition}[Schedule]\label{def:schedule}
A schedule for $(\lambda,\eta,\xi)$ starts at $\lambda$, ends at a
point of $W\xi$, and, between finitely many ordered switch times,
satisfies $y'=w\eta'$ for a fixed $w\in W$.  At a switch at time $t$,
the frame changes from $w$ to $s_\beta w$, where $\beta$ is a signed
root satisfying
\begin{equation}\label{real:signs}
 G=y(t)\cdot\beta\ge0,
 \qquad -\beta\cdot w\mu\ge0.
\end{equation}
The number $G$ is the position gap, and $-\beta\cdot w\mu$ is the
total moving-weight pairing.  Several switches may occur at the same
time, with their order part of the data.  The switch coefficient is
$m=-\beta\cdot w\eta(t)$.  The schedule is \emph{integral} if every
switch coefficient is an integer.
\end{definition}

If the current frame is $w$ and the preceding switches have coefficients
$m_i$ and position roots $\beta_i$, continuity gives
\begin{equation}\label{real:balance}
 y(t)=\lambda+w\eta(t)-\sum_{i\text{ preceding }t}m_i\beta_i.
\end{equation}
The formula uses the chosen ordered prefix at a multiple switch time.
Indeed $s_\beta w\eta=w\eta+m\beta$ at a switch, exactly canceling
the newly added term in the sum.

There are three useful simplifications.  A switch with
$-\beta\cdot w\mu=0$ changes no shape increment: the root
$w^{-1}\beta$ has one sign on the dominant chamber, so its zero total
pairing forces zero pairing with every increment.  Such a switch can
be deleted, changing later direction indices on the right by a
stabilizer of the entire shape if necessary.  Switches at $0$ or $1$
can be absorbed into the initial frame or the endpoint direction.
Finally, if $G=0$, reflect the whole suffix, including its switch
roots, in that wall and delete the switch.  The path remains continuous
and its endpoint remains in $W\xi$.  These operations also preserve
integrality when the original schedule is integral.

Let $\mathcal G$ be the simply connected complex semisimple group
and $K$ a compact form; the matrix models below describe their compact
orbit quotients.  We identify real Cartan vectors with Hermitian generators in
$i\mathfrak k$.  A \emph{real moment solution} for
$(\lambda,\mu,\xi)$ is a zero-sum triple in the compact orbits of
$\lambda,\mu,-\xi$.  In type $D$, these are $i$ times real
skew-symmetric matrices.  Their eigenvalues are the positive and
negative signed coordinates; the compact orbits are special
orthogonal orbits, so their orientation information is retained.  In
permutation systems we use Hermitian matrices.  The central equation
is imposed separately.  Compactness of $K$ shows that moment
solvability is closed in the three labels.

Write $B=T U^+$ and $B^-=T U^-$, with $B$ fixing highest lines.
We index the torus-fixed points of $\mathcal G/B$ by $W$, and use the
Bruhat convention
\[
 u\ge v\quad\Longleftrightarrow\quad
 vB\in\overline{BuB/B}.
\]
Downward reflection steps have the form $u\mapsto s_\beta u$ with
$\beta>0$ and $u^{-1}\beta<0$.  A Bruhat cover is a strict
comparison with no intervening Weyl element.  Inversion preserves
Bruhat order.
For a regular dominant $\rho$, the flow
$a(t)=\exp(t\rho)$ has past set $U^+p$ and future set $U^-p$ at a
fixed flag $p$.

The principal construction has the following precise output.  Its
cover condition is read in the chamber coordinates obtained by
folding the path into $\mathcal C$.

\begin{proposition}[Basic real schedule]\label{prop:basic-schedule}
Let the root system be of type $D$ or a permutation system, possibly
with an orthogonal central space.  If $(\lambda,\mu,\xi)$ has a real
moment solution, then every shape $\eta$ ending at $\mu$ admits a
schedule whose switch roots form a real basis of the semisimple
space.  Weak signs, endpoint times, and coincident times are allowed.
More precisely, it is a limit of schedules for nearby regular data
whose switches are covers in their local chamber frames.  Simplifying
the limiting schedule may leave fewer than rank-many switches.
\end{proposition}

We prove the proposition after constructing decreasing paths and
maximizing their area.  First we develop the coefficient tools needed
for the converse from schedules to moments.

\subsection{Coefficients and flags}

The two tools below also combine the fundamental columns in
Section~\ref{sec:columns}.  The classical orthogonal representation
models are recalled in \cite[Lectures~19--20]{FH91}.  We give the
coefficient and composition arguments explicitly because the second
must control restriction to a Cartan summand.

For an integral dominant $h$, write $v_h$ for a highest vector of
$V(h)$, and write $f_i$ for the negative simple generators.  A flag
vector of weight $h$ at a flag $b$ means a nonzero vector on the
translate of the line $\C v_h$ associated with $b$.

The coefficient description below is the dual form of the
Parthasarathy--Ranga Rao--Varadarajan tensor-multiplicity formula
\cite{PRV67}; see also \cite[Theorems~1.1--1.2]{PY08}.  For the classical
coefficient criterion, see Kumar~\cite[Theorem~3.7]{Kumar2010}, where it
is attributed to Kostant. The already-dominant extremal consequence
appears in \cite[Corollary~3.8]{Kumar2010}. We include the argument to
fix the highest-coordinate convention.

\begin{lemma}[Highest-to-highest coefficients]\label{lem:coefficient}
Let $h,h',\theta$ be dominant integral weights.  The linear forms
obtained by taking the highest-output coordinate of
\[
 T(v_h\otimes {\cdot}),\qquad
 T:V(h)\otimes V(\theta)\longrightarrow V(h'),
\]
are exactly the forms supported at weight $h'-h$ that annihilate
\begin{equation}\label{real:coefficient-images}
 \sum_i f_i^{h\cdot\alpha_i+1}V(\theta).
\end{equation}
In particular, if $h'=h+u\theta$ is dominant, the extremal-coordinate
functional at $u\theta$ gives such an intertwiner and is nonzero on
the flag vector at $u$.
\end{lemma}

\begin{proof}
Take coinvariants for the diagonal negative nilpotent algebra on the
tensor product.  At weight $h'$ this quotient records exactly the
highest coordinate of each copy of $V(h')$.  Move negative generators
from the first tensor factor to the second by the antipode.  The
highest-generator presentation
\[
 V(h)=U(\mathfrak n^-)
 \Big/\sum_i U(\mathfrak n^-)f_i^{h\cdot\alpha_i+1}
\]
then gives precisely the quotient by
\eqref{real:coefficient-images}: the antipode reverses products, so a
relation $Xf_i^{h\cdot\alpha_i+1}$ transfers into the image of the
indicated power on an arbitrary second vector.  This proves both
necessity and sufficiency after dualizing the weight quotient.

The presentation itself follows by imposing the simple singular-vector
relations on the Verma module.  Each simple $\mathfrak{sl}_2$ then
acts integrably: local nilpotence extends from the generating vector
using the adjoint action.  The weights are therefore Weyl-invariant
and bounded in the highest-weight order; together with PBW this makes
the quotient finite-dimensional.  Complete reducibility identifies
it with the irreducible highest-weight module.

For an extremal weight $u\theta$, the upward simple-string length is
$\max(0,-u\theta\cdot\alpha_i)$.  It is at most
$h\cdot\alpha_i$ because $h+u\theta$ is dominant.  Thus its
coordinate functional annihilates every image in
\eqref{real:coefficient-images}.
\end{proof}

Products of suitably ordered flag coordinates are central to standard
monomial theory; see, for example, \cite[Theorem~4]{Lit98}. The
following form works directly with the coefficients just constructed.
Its ordering hypothesis is what ensures that a nonzero product survives
on the Cartan summand.

\begin{lemma}[Flag composition]\label{lem:flag-composition}
Suppose that $h_0,\ldots,h_s$ and $\theta_1,\ldots,\theta_s$ are
dominant integral weights and that there are intertwiners
\[
 T_j:V(h_{j-1})\otimes V(\theta_j)\longrightarrow V(h_j).
\]
Assume their highest-to-highest coefficients are nonzero on flag
vectors at flags $b_j$.  If the past and future limits of $b_j$ under
$a(t)$ are $u_j,v_j$, and $v_j\ge u_{j+1}$, then the iterated
composition is nonzero on
\[
 V(h_0)\otimes V(\theta_1+\cdots+\theta_s).
\]
Here the second factor is the Cartan summand in the tensor product of
the $V(\theta_j)$.
\end{lemma}

\begin{proof}
The empty composition is the identity; assume $s\ge1$.
We first explain why one positive-flow orbit can approximate the
waypoints $b_1,\ldots,b_s$ in their prescribed order, with successive
time gaps tending to infinity.  When $v\ge u$, there is a bridge in
$U^+v\cap U^-u$.  Indeed $U^+v$ approaches $u$, and a neighborhood
of $u$ is covered by small $U^+U^-$ translates.  Removing the small
positive factor from a point of $U^+v$ leaves a point in that
intersection.

It remains to glue at a shared fixed flag $p$.  Suppose the preceding
waypoint approximation lies at $n^-p\in U^-p$ and the next desired
point is $n^+p\in U^+p$.  Replace the former by
\begin{equation}\label{real:gluing-point}
 n^-a(-t)n^+a(t)p.
\end{equation}
As $t\to\infty$ this tends to $n^-p$, whereas after flowing for time
$t$ it tends to $n^+p$.  The remaining negative factor after that
flow preserves the future cell containing $n^+p$.  Insert the bridges
and repeat this construction, choosing each new time gap sufficiently
large to preserve every earlier approximation.

For each such approximation, choose flag vectors $z_j\in V(\theta_j)$
at its common starting flag $z$, and times $t_1<\cdots<t_s$.
Let $C$ denote the iterated composition of the $T_j$, and set
\[
 F_j(x)=T_j\bigl(x\otimes a(t_j)z_j\bigr),\qquad
 C_z(x)=C(x\otimes z_1\otimes\cdots\otimes z_s).
\]
Equivariance gives the exact identity
\[
 F_s a(t_s-t_{s-1})F_{s-1}\cdots a(t_2-t_1)F_1
   =a(t_s)C_z a(-t_1).
\]
Indeed, writing each $F_j$ as
$a(t_j)T_j(a(-t_j){\cdot}\otimes z_j)$ cancels all the
intermediate frame factors.  For $s=1$ the identity is just this
equivariance formula.

Normalize each translated flag input so that it approaches the
prescribed vector at $b_j$.  Normalize each intervening
$a(t_{j+1}-t_j)$ by its highest eigenvalue on $V(h_j)$; as the time
gaps tend to infinity, it approaches projection onto the highest
line.  The highest-input, highest-output coordinate of the left side
therefore tends to the product of the nonzero waypoint coefficients.
It is nonzero for some finite approximation.  On the right side the
outer frame factors act by nonzero scalars on those coordinates, so
$C_z$ is nonzero.  The $z_j$ lie at one common flag, and hence their
tensor belongs to the Cartan summand.  This proves the assertion.
\end{proof}

The next compact maximization argument is a direct instance of the
norm methods of Kempf and Ness \cite{KN79}. We need only the
invariant-to-moment implication, whose proof follows.

\begin{lemma}[Invariants give moments]\label{lem:invariant-moment}
A nonzero invariant in
$V(\lambda)\otimes V(\mu)\otimes V(\xi^*)$ gives a real moment
solution for $(\lambda,\mu,\xi)$.  The same conclusion follows from
an invariant for simultaneous positive integer multiples of the labels.
\end{lemma}

\begin{proof}
Take an invariant functional and maximize its absolute evaluation on
triples of unit flag vectors, using compact-invariant Hermitian norms.
The maximum is attained and is nonzero, since the flag-vector orbits
span their representations.  Vary all three vectors simultaneously
by a Hermitian generator $H$ and renormalize.  Invariance makes the
unnormalized numerator constant; differentiation of the three norms
therefore gives zero sum of the expectations of $H$.  Since this
holds for all $H$, the three moments sum to zero.  The moment of a
compact translate of a unit highest vector is the corresponding
compact translate of its weight.  This is the required solution.
For scaled labels divide the three Hermitian elements by the scale.
\end{proof}

\begin{lemma}[Schedules give moments]\label{lem:schedule-moment}
A real schedule for $(\lambda,\eta,\xi)$ gives a real moment
solution for $(\lambda,\mu,\xi)$, where $\mu=\eta(1)$.
No lattice condition is required.  The statement also holds for
root subsystems with arbitrary orthogonal central coordinates.
\end{lemma}

\begin{proof}
Apply the three schedule simplifications.  The remaining switch gaps,
opposite moving-weight pairings, and distances from the time endpoints
are strict.  Perturb the initial position, shape, and times to rational
generic data, allowing the shape endpoint and the final position to
vary.  Make the shape slopes regular, separate the switch times, and
make all ordinary wall crossings simple, transversal, and disjoint
from the switches and shape breaks.  All strict signs persist.

Project the path to the dominant chamber.  At an ordinary crossing,
the incoming folded velocity has negative pairing with the crossed
simple root, so the direction makes a downward simple Bruhat step.
At an original switch, its positive gap makes its root positive in
the local chamber frame, and the opposite pairing with the regular
moving weight makes its inverse root negative.  It too is a downward
Bruhat reflection.  Shape breaks alone do not change the Weyl frame.

Subdivide at every crossing, switch, and shape break.  A common
integer dilation makes the semisimple positions and shape increments
integral.  Each segment then has an extremal intertwiner from
Lemma~\ref{lem:coefficient}; its waypoint is the fixed flag of the
segment's direction.  Their decreasing order permits
Lemma~\ref{lem:flag-composition}.  The resulting nonzero intertwiner
gives a real moment solution by Lemma~\ref{lem:invariant-moment}.
Divide by the dilation and pass to the unperturbed labels by
compact-orbit closedness.  Central coordinates are omitted from the
representation argument and restored by their unchanged addition
equation; no integrality of a central coordinate is needed.
\end{proof}

\subsection{From moments to decreasing real paths}

We next establish the real paths on which the area argument operates.
During this construction we may freely perturb all three labels before
returning to their prescribed values.

For a full flag $F$ and a fundamental weight $\omega$, let
$m_\omega(F,H)$ be the least eigenvalue of a Hermitian generator $H$
occurring in a flag vector of weight $\omega$ at $F$.  Extend this
additively to dominant real weights.  A triple of weighted flags
$(F_j,\theta_j)$ is \emph{strict} if
\begin{equation}\label{real:strict}
 \sum_{j=1}^3 m_{\theta_j}(F_j,H)<0
 \qquad (H\in i\mathfrak k,\ |H|=1).
\end{equation}
Replacing $<$ by $\le$ defines semistability.  We work in the
semisimple space in this definition.

\begin{lemma}[Strict approximation]\label{real:strict-approximation}
Every real moment triple can be approximated by triples of regular
labels admitting strict flags.  Such strict label triples form an
open set.  In particular, after choosing a nearby regular shape,
there is open freedom to vary its two endpoint labels independently.
\end{lemma}

\begin{proof}
A moment solution supplies semistable flags: each least supported
level is at most its expectation, and the three expectations sum
to zero.

For fixed labels, the strict and semistable loci of triples of flags
are Zariski open.  To see the finite incidence argument, write a test
as $H=\Ad(k)d$ with $d$ dominant.  A lower bound on a least level
means that each relevant highest line lies in the $k$-translate of an
upper level space
\[
 \bigoplus_{\chi\cdot d\ge c} V_\chi.
\]
Such spaces are $B$-stable, and only finitely many choices occur in
the finitely many fundamental representations.  For each choice,
incidence with a common translating flag is closed; its projection
is closed because the flag variety is projective.  Among these
finitely many incidence types, retain those for which some dominant
$d$ violates the prescribed strict or weak inequality.  Their union
is exactly the failure locus.

Strict regular data exist.  Take regular Hermitian elements $X,Y$
with zero common semisimple Lie centralizer and with $-X-Y$ regular.
For example, one can start with $X$ regular toral and choose $Y$
generically with nonzero components on every simple root.  Attach
the three compact flags realizing these representatives.  Equality
in the sum of least-level inequalities forces each flag vector to
be an eigenvector of the testing generator.  For regular weights
this makes the generator commute with all three representatives,
so it must be zero.  The resulting triple is strict.

Intersect the nonempty open flag locus for the original semistable
labels with that for this auxiliary strict triple.  The product of
flag varieties is irreducible, so the intersection is nonempty.
At a common triple of flags, a positive convex combination of the
original and auxiliary labels is strict and regular.  Such labels
approach the original labels.  Finally, the least supported level
is upper semicontinuous in $H$: disappearing components can only
raise its value at the limit.  Compactness of the unit sphere makes
the strict margin uniform.  Small changes of the finitely many
weight coefficients preserve that margin.  This proves openness
and the stated freedom of endpoints.
\end{proof}

The relation between weighted configurations and symmetric-space
polygons was developed by Kapovich, Leeb, and Millson
\cite[Section~5.3, Proposition~5.6]{KLM09}. We give the strict three-point construction
needed here, including the inward estimate for the composition of
pushes.

\begin{lemma}[Symmetric-space triangles]\label{real:triangles}
If $(\lambda,\mu,\xi^*)$ admits strict flags, then for every $R>0$
there is a triangle in $\mathcal G/K$ with successive dominant side
vectors $R\lambda,R\mu,R\xi^*$.
\end{lemma}

\begin{proof}
The push of length $R\theta$ from $hK$ toward a full flag $F$ is the
path
\[
 hk\exp(s\theta)K,\qquad 0\le s\le R,
 \qquad h^{-1}F=kB,\quad k\in K.
\]
It is continuous and well-defined, and the same flag specifies its
forward direction all along the push.  Compose the three pushes.
We prove that this composition has a fixed point.

In polar coordinates $\exp(P)K$, radial distance is $|P|$, and inverse
polar coordinates are distance non-increasing.  One direct verification
uses the differential of the exponential: after translation back by
$\exp(-P)$, its Hermitian part is
\[
 \frac{\sinh(\operatorname{ad}P)}{\operatorname{ad}P}\,\delta P.
\]
This expands lengths and is the identity in the radial direction.
Polar decomposition gives global coordinates.  In particular, points
at a bounded mutual distance and tending to infinity have the same
limiting radial unit direction.

At $\exp(lH)K$, $|H|=1$, the radial derivative of a push toward $F$
is the pairing of $H$ with the compact moment of $\theta$ at
$\exp(-lH)F$.  If $l\to\infty$ and $H\to H_0$, then
\begin{equation}\label{real:radial-bound}
 \limsup \frac{\dd}{\dd s}|P(s)|
 \le m_\theta(F,H_0).
\end{equation}
For a fundamental representation, take a narrow spectral window
around the least supported level for $H_0$.  Before applying the
exponential, the flag vector's projection onto a slightly enlarged
window has norm bounded below for all nearby $H$.  Multiplication
by $\exp(-lH)$ and normalization suppresses levels above that window
exponentially.  The resulting expectation has limsup at most the
window's upper endpoint.  Shrinking the window proves
\eqref{real:radial-bound}; addition proves it for real dominant
weights.

For fixed $R$, all three pushes have bounded displacement.  Along
any sequence tending to infinity with limiting unit direction $H_0$,
the integrated radial derivatives therefore have total limsup at most
\[
 R\sum_{j=1}^3 m_{\theta_j}(F_j,H_0)<0.
\]
The derivatives are bounded, so this estimate passes through the
finite integrals.  It follows, by compactness of unit directions,
that every sufficiently large sphere maps strictly inward under the
composition.  Compose with radial retraction onto the corresponding
closed ball and apply Brouwer's fixed-point theorem \cite{Brouwer1911}.
A fixed point
cannot lie on its boundary, by the strict inward property; an interior
fixed point of the retracted map is a fixed point of the original
composition.  The three pushes from that point form the triangle.
\end{proof}

\begin{lemma}[Decreasing radial paths]\label{real:decreasing-paths}
For generic strict labels and a generic regular polygonal shape
ending at $\mu$, there is a dominant path from $\lambda$ to $\xi$
whose velocities are Weyl translates of the shape velocities, in
weakly Bruhat-decreasing order.
\end{lemma}

\begin{proof}
Move the first vertex of the triangle in Lemma~\ref{real:triangles}
to $K$.  Replace its middle side by
\[
 g_R\exp(R\eta(t))K.
\]
Here $g_R$ includes the entering direction frame, so the replacement
has the same endpoints.  Its initial and final radial Cartan parts
are $R\lambda$ and $R\xi$, the latter because the last side is
reversed.  Denote its radial Cartan part divided by $R$ by $x_R(t)$.

Let $R\to\infty$.  Use all fundamental representations and a regular integral
highest-weight representation $V(\sigma)$, with compact-invariant
norms.  We can choose a common flag $z_R$ such that, in each of these
representations and for every nonzero Cartan-weight projection,
\begin{equation}\label{real:operator-comparison}
 \|\operatorname{pr}_\chi g_R^\dagger z_R\|
 \ge c\,\|\operatorname{pr}_\chi g_R^\dagger\|,
 \qquad c>0
\end{equation}
with $c$ independent of $R$; use unit flag vectors for $z_R$.
Indeed normalize the finitely many nonzero projected operators to
norm one.  There are only finitely many possible nonzero-index patterns;
for each pattern their tuples range in a compact space.  A generic flag
vector avoids the kernels of any such tuple, because each nonzero
operator is nonzero on some flag vector and a finite union of proper
closed subsets does not fill the flag variety.  The maximum, over
flags, of the minimum of the finitely many norms is positive and
continuous in the tuple.  Its compact minimum gives $c$.

The adjoint operation is compatible across these representations.
The highest-weight operator norm formula from compact-Cartan
decomposition, applied to the adjoint product, and
\eqref{real:operator-comparison} show that, after a subsequence,
\begin{equation}\label{real:tropical}
 \theta\cdot x(t)
 =\max_\chi\bigl(c_{\theta,\chi}+\chi\cdot\eta(t)\bigr).
\end{equation}
Here $c_{\theta,\chi}$ is the limit of
$R^{-1}\log\|\operatorname{pr}_\chi g_R^\dagger z_R\|$;
values $-\infty$ are allowed.  To obtain the formula, the squared norm of
$\exp(R\eta(t))g_R^\dagger z_R$ is the sum of the squared norms
of its weight projections.  The logarithm of its norm therefore
differs from the largest logarithm of a projection norm by a bounded
amount.  The logarithmic error in
\eqref{real:operator-comparison} is also bounded before division by
$R$.  The coefficients are bounded above by the operator norm of
$g_R^\dagger$, whose Cartan part is $R\lambda$.  Thus subsequences
exist simultaneously for all finitely many coefficients.  In each
fundamental representation, operator-norm submultiplicativity applied
to $\exp(R(\eta(t)-\eta(s)))$ and its inverse bounds
$|\theta\cdot(x_R(t)-x_R(s))|$ by
$\max_\chi |\chi\cdot(\eta(t)-\eta(s))|$.  The fundamental
weights span the Cartan space, so the paths $x_R$ are
uniformly Lipschitz.  Their fixed endpoints make the limit a dominant
path from $\lambda$ to $\xi$.

Choose the slopes of $\eta$ to separate all distinct weight slopes
of $V(\sigma)$.  Except at finitely many times, the maximum in its
version of \eqref{real:tropical} has a unique weight.  Every limit
of the flags $\exp(R\eta(t))g_R^\dagger z_R$ then lies in that
weight space in the regular highest-weight embedding.  Such a flag
is torus-fixed.  The fixed flags have distinct extremal weights,
so it is a single fixed flag $p$ on each open interval.  The
fundamental formulas give $x'=p^{-1}\eta'$ there.

If a later interval gives $q$, then $p\ge q$.  To see this, use a
local factorization near $p$ as small $U^-U^+$ translates of $p$.
Multiplication by the exponential of the scaled regular dominant
increment contracts the conjugated $U^-$ factor to the identity.
On the compact flag variety this factor acts uniformly close to the
identity, even if the remaining $U^+$ factor diverges.
Consequently $q$ belongs to the closure of $U^+p$, which is exactly
the asserted Bruhat comparison.  Inversion preserves Bruhat order,
so the velocity frames $p^{-1}$ are weakly decreasing as well.
\end{proof}

\subsection{Area maximization and the basis of switches}

We now refine the decreasing path so that its nonzero changes are
single root reflections and their roots form a basis. Area maximization
first controls linear dependencies among changes of velocity. The
flag-orbit polytope argument then identifies those changes as reflections
and proves the cover property needed in the comparison count.

Fix a regular dominant vector $\rho$ and write
\[
 \mathcal A(x)=\int_0^1\rho\cdot x(t)\,\dd t.
\]
For a fixed regular shape and fixed endpoints, consider the dominant
paths supplied by Lemma~\ref{real:decreasing-paths}, allowing all
weakly Bruhat-decreasing direction sequences.  This is a compact
set: there are finitely many strict direction chains, their ordered
time simplices are compact, and the endpoint and chamber conditions
are closed.  Hence $\mathcal A$ has a maximum.

We shall use the following parabolic consequence of the subword
criterion to compare directions when a path meets chamber walls.
If $W_I$ is a standard parabolic subgroup, replacing an element by
the maximal representative of its left $W_I$-coset preserves Bruhat
comparisons.  Moreover,
\begin{equation}\label{real:parabolic}
 u\ge v,\quad u,v\text{ left }I\text{-maximal}
 \quad\Longrightarrow\quad gu\ge gv\quad(g\in W_I).
\end{equation}
For completeness, projection to the minimal coset representative
preserves comparisons: in a reduced factorization with parabolic
prefix, the subword criterion gives the comparison of the remaining
minimal representatives.  Adding the common length-additive prefix
$w_I$ gives the first assertion.  Writing
$u=w_I\bar u$, $v=w_I\bar v$ and adding the common
length-additive prefix $gw_I$ gives \eqref{real:parabolic}.
Here $w_I$ denotes the longest element of $W_I$.

At a direction jump $u\to v$ of such a path, let $I$ consist of the
simple walls containing its position.  Regularity of the adjacent
slopes makes $u$ left $I$-maximal and $v$ left $I$-minimal.  Reflect
the suffix by $w_I$.  The jump is now
\begin{equation}\label{real:unfolded-jump}
 u\longrightarrow w_Iv,\qquad u\ge w_Iv,
\end{equation}
and both states are left $I$-maximal.  Perform these operations in
order, carrying each subsequent instruction by the accumulated
left frame.  We call this \emph{unfolding}.  The original path is
recovered by projection to the dominant chamber.  An unfolded jump
is \emph{nonzero} if its two states differ; otherwise it disappears.
The unfolded terminal point is a fixed Weyl translate of $\xi$.
Endpoint displacements will be computed in this common unfolded
space, whereas the area is always that of the folded dominant path.

\begin{lemma}[Area variation]\label{real:variation}
For an area-maximizing path with regular shape, the velocity differences
at its nonzero unfolded jumps away from shape breaks are linearly
independent.  Here all differences are measured in the common unfolded
coordinate space.
\end{lemma}

\begin{proof}
\emph{An endpoint-preserving variation.}
Suppose those differences have a nontrivial linear dependence.  Move
the corresponding jump times linearly in a parameter $s$ in the
dependency direction, leaving the shape-break times fixed.  For small
positive and negative $s$, the moved times stay in their original
slope pieces.  The unfolded endpoint change is exactly the sum of
the velocity differences times the time changes, so it is zero.

\emph{Admissibility after folding.}
We check that folding these variations still gives admissible
Bruhat-decreasing paths.  Near an old jump position, only the wall
subgroup $W_I$ can be needed for sorting.  Both states in
\eqref{real:unfolded-jump} are left $I$-maximal.  Consequently
\eqref{real:parabolic} preserves their comparison after any common
local sorter $g\in W_I$.  Ordinary crossings introduced on either
side are downward simple steps.  If a moved jump and a wall crossing
are simultaneous, sort on the incoming side, make the comparison,
then perform the pure crossing on the outgoing direction.  To verify
the last step, the positive roots of the approached walls pull back
under the incoming sorter to positive roots of the local $I$-system:
the incoming slope has negative pairing with them.  The other
$I$-maximal state is also negative-facing on that approached
subsystem.  The outgoing pure sorting is therefore downward.  The
same reasoning works when a fixed shape break shares the old site.
Away from the old sites and breaks there are no new wall contacts;
a regular slope cannot run along a wall of a dominant path.

\emph{First and second variations of the folded area.}
We first show that $\mathcal A(x_s)$ has a common two-sided derivative
at zero.  The issue is the change of chamber sorter near old wall
sites.  There
are finitely many jump, break, and wall-contact times.  On each
adjacent open segment, regularity of the slope makes every vanishing
wall coordinate grow at least linearly with distance from its contact
time.  The unfolded perturbation is uniformly $O(|s|)$.  Thus, outside
neighborhoods of radius $C|s|$ of those finitely many times, with $C$
fixed sufficiently large, both signs of the perturbation have the
same chamber sorter as the original path.  They also have the same
ordered prefix of moved jumps there.  Consequently their unfolded
positions have the form $y_s(t)=y_0(t)+s z(t)$, with the same bounded
$z(t)$ for both parameter signs, and chamber projection applies the
same linear Weyl frame.  Chamber projection is globally Lipschitz
(it is a continuous, finite piecewise-orthogonal map), so the omitted
set, of length $O(|s|)$, contributes only $O(s^2)$ to the area
difference.  Omitting it from the bounded first-variation integrand
also changes that integral by $O(|s|)$.  There is therefore a common
two-sided first variation, which must vanish at an area maximum.

On either sufficiently small one-sided parameter interval, all
resulting sites, including sorting crossings, have affine times
$\tau_j(s)$.  Indeed positions are affine in $(t,s)$ in each fixed
frame, and crossing times solve linear wall equations.  Write
$v_j^-,v_j^+$ for the actual folded velocity vectors immediately
before and after the $j$th site.  After fixing the one-sided order
of these finitely many sites, integration of the piecewise constant
velocities with weight $1-t$ gives
\begin{equation}\label{real:second-variation}
 \frac{\dd^2}{\dd s^2}\mathcal A(x_s)
 =\sum_j \rho\cdot(v_j^+-v_j^-)
                 \bigl(\tau_j'(s)\bigr)^2.
\end{equation}
A fixed shape break contributes zero.  Every other term is
nonnegative, because the corresponding comparison is downward and
the local shape slope is regular dominant.  A moved nonzero original
jump gives a strictly positive term: its two unfolded states lie in
distinct $W_I$-cosets, so sorting cannot erase their difference.
The first variation vanishes and the one-sided second variation is
strictly positive, contradicting maximality.
\end{proof}

Choose the endpoints generically in the open set supplied by
Lemma~\ref{real:strict-approximation}, with the regular shape fixed.
The independent differences in Lemma~\ref{real:variation} then form
a basis of the semisimple space.  Indeed, for any fixed discrete
unfolding data, the endpoint balance has the form
\begin{equation}\label{real:time-balance}
 y(1)-y(0)=b+\sum_j \tau_j(a_j-b_j),
\end{equation}
where the summation is over the nonzero nonbreak jumps, and the
fixed vector $b$ includes every contribution from the shape breaks.
There are only finitely many discrete data.  A deficient span in
\eqref{real:time-balance} confines the freely varying ends to a
proper affine subspace, which generic ends avoid.  There are therefore
exactly as many such jumps as the semisimple rank, and their times
are determined by the balance.

For the same generic choice, every nonzero jump and every shape break
lies off all walls in its local frame.  For fixed discrete data,
translate both unfolded ends by the same small vector.  The balance
and the basis times stay fixed, while all their positions translate
by that vector.  Each asserted wall equality is thus a proper affine
condition on the free endpoints.  Avoid the finitely many such
conditions.  Pure sorting crossings are not nonzero unfolded jumps
and impose no extra restriction here.

\subsection{From a basis to cover reflections}

The nonbreak jump differences now form a basis, and every nonzero
jump site is off the chamber walls.  The basis allows us to insert a
new direction at any such site and compensate its endpoint displacement
by shifting the existing jump times.  We use these variations to
expose an edge, and then to show that its reflection is a cover.

\begin{lemma}[Strict exposure]\label{real:exposure}
For such a generic area maximum, let $u\to v$ be a nonzero jump in
local chamber coordinates.
\begin{enumerate}[label=\textup{(\roman*)}]
\item If the site is not a shape break and the slope is $d$, there
is a linear functional $f$ whose maximum among
$\{hd:v\le h\le u\}$ is attained exactly at $ud$ and $vd$.
\item If it is a shape break with slopes $d_-,d_+$, a common
functional $f$ has strict maximum at $ud_-$ among the corresponding
$d_-$-vectors, and strict maximum at $vd_+$ among the corresponding
$d_+$-vectors.
\end{enumerate}
\end{lemma}

\begin{proof}
The derivative of the unfolded endpoint with respect to the nonbreak
jump times is an isomorphism: its columns are the basis of velocity
differences.  We can therefore add a linear functional of the
unfolded terminal position to the area objective so that all
unconstrained basis-time derivatives vanish.  This functional is
constant on the endpoint constraint, so the constrained maximization
problem is unchanged.

At the chosen site, insert an intermediate state $h$ from the Bruhat
interval for a small duration, on the left or on the right.  Its
endpoint displacement is exactly compensated by linear shifts of
the basis times.  For sufficiently small duration, all these times
remain in their fixed slope pieces, so compensation is exact, not
merely to first order.  If the chosen nonbreak site itself supplies
a compensating variable, shift the inserted group as a whole.
Chamber feasibility follows from the sorting argument in
Lemma~\ref{real:variation}.

Use the fixed Weyl frame at the chosen off-wall site to identify
its local coordinates with the common unfolded space.  The first
variation of the modified objective is then a linear functional $f$
of the inserted local velocity difference.  It is the same
from either side of the site: it integrates the effect of that
displacement on the later folded path, together with the terminal
multiplier, and changing the lower time endpoint by a vanishing
interval does not change the first variation.  At a nonbreak site,
shifting the original jump has vanishing first variation, so
$f(ud)=f(vd)$.  Maximality gives $f(hd)\le f(ud)$.
At a break, a left insertion gives
$f(hd_-)\le f(ud_-)$ and a right insertion gives
$f(hd_+)\le f(vd_+)$.  These include extension of the opposite
endpoint state across the break.

For an inserted interval parameterized by its duration, the two
boundary-time derivatives differ by one, so compensation cannot fix
both.  Thus each genuine intermediate insertion, or opposite-state
extension at a break, moves at least one nonzero jump time.  If its first-order
inequality were equality, Equation~\eqref{real:second-variation}
would give a strictly positive one-sided second variation and hence
an improvement.  Therefore every such inequality is strict.  This
proves both assertions.
\end{proof}

The orbit-closure polytope and its root-parallel edges are classical:
see Atiyah \cite[Theorem~2]{Atiyah82} and Gel'fand--Serganova
\cite[Section~7.1, Theorem~1]{GeS87}. We include the proof in the form
needed here, keeping track also of the fixed-flag labels as a regular
real highest weight varies.

\begin{lemma}[Flag orbit polytopes]\label{real:polytope}
Choose a flag $b\in U^+u\cap U^-v$, with $u\ge v$.  The torus-fixed
points of its torus-orbit closure lie in $[v,u]$ and include $u,v$.
For a regular dominant real weight $d$, take the convex hull of the
corresponding vectors $hd$.  Its normal fan, with vertices labeled
by fixed flags, is independent of $d$.  Every edge joins two vertices
related by a root reflection.
\end{lemma}

\begin{proof}
The orbit closure is contained in the Schubert closure through $u$
and the opposite Schubert closure through $v$, giving the interval
assertion; the two flow limits supply its endpoints.  For integral
regular $d$, this is the weight polytope of the torus orbit of the
line of the flag vector in $V(d)$.  Vertices are obtained by generic
one-parameter limits, and the fixed flags in this embedding have
Weyl-extremal weights.

The flag vector for a sum of fundamental weights is the tensor of
their flag vectors in the Cartan summand.  Its torus weight polytope
is therefore the corresponding Minkowski sum of the fundamental
polytopes.  For strictly positive coefficients its normal fan is
their common refinement, independent of those coefficients.  The
labels are independent as well, since a generic torus limit is the
same flag in all highest-line maps.  Positive real coefficients give
the same conclusion for real regular $d$.

For integral $d$, an edge is exposed by a one-parameter limit and
comes from a one-dimensional torus orbit in the flag variety.  Its
codimension-one stabilizing subtorus annihilates a tangent root at a
fixed endpoint, since the tangent weights of the flag variety there
are roots.  Thus the edge direction is a root direction.  Its two
Weyl vertices have equal norm, and two distinct equal-norm points on
a line in root direction are exchanged by the corresponding root
reflection.  The common labeled fan carries this statement to every
regular real $d$.
\end{proof}

\begin{lemma}[Exposed jumps are covers]\label{real:covers}
The generic area maximum has no nonzero jump at a shape break.
Every other nonzero unfolded jump is a reflection and is a Bruhat
cover in its local chamber frame.
\end{lemma}

\begin{proof}
Choose the flag in Lemma~\ref{real:polytope}.  At a shape break,
Lemma~\ref{real:exposure} would give two different unique maximizing
vertex labels for the same functional on two polytopes with the
same labeled normal fan.  This is impossible.  At a nonbreak site,
that lemma exposes exactly the endpoint vertices, so they form an
edge.  Lemma~\ref{real:polytope} makes the jump a reflection.

Write it as $u\to s_\beta u$, with $\beta>0$ and
$u^{-1}\beta<0$.  Suppose it is not a Bruhat cover.  In these simply
laced systems there are positive roots $\alpha,\delta$ such that
\begin{equation}\label{real:inversion-pair}
 \beta=\alpha+\delta,\qquad
 u^{-1}\alpha<0,\quad u^{-1}\delta<0.
\end{equation}
Here is the inversion count proving the assertion.  Apart from
$\beta$, the positive roots taken negative by $s_\beta$ are paired
as $\alpha,\beta-\alpha$.  A pair with both roots inverted by
$u^{-1}$ contributes two to the length drop, whereas a mixed pair
contributes zero.  A pair with neither root inverted is impossible,
since their sum has negative inverse image $u^{-1}\beta$.
Thus a drop greater than one forces \eqref{real:inversion-pair}.

In the rank-two subsystem on these three positive roots, both
$s_\alpha u$ and $s_\delta u$ are proper intermediate states between
$u$ and $s_\beta u$.  Strict exposure therefore forces
$f\cdot\alpha<0$ and $f\cdot\delta<0$: for example,
$s_\alpha ud-ud=-(ud\cdot\alpha)\alpha$ is a positive multiple
of $\alpha$.  Equality of the two exposed endpoint values gives
$f\cdot\beta=0$, contrary to \eqref{real:inversion-pair}.
The reflection is a cover.
\end{proof}

\begin{proof}[Proof of Proposition~\ref{prop:basic-schedule}]
First use generic strict labels and a generic regular shape.
Lemma~\ref{real:decreasing-paths} supplies the nonempty compact path
class.  At an area maximum, Lemma~\ref{real:variation} and the generic
endpoint balance give a basis of nonbreak velocity differences;
Lemma~\ref{real:covers} removes all nonzero break jumps and makes
every remaining nonzero jump a cover reflection.  Its velocity
difference is a nonzero scalar multiple of its root.  Thus those
roots form a basis.  In the unfolded coordinates they obey
\eqref{real:signs}: the local comparison root has nonnegative
position pairing and negative pairing with the incoming moving
weight.  They therefore give a schedule.

For the prescribed moment data, use Lemma~\ref{real:strict-approximation}
and regular shapes converging to the prescribed shape.  Such shapes
can be obtained by small positive interpolation with shapes ending at
the perturbing regular weights, followed by generic choices of their
slopes and break values.  At each stage the strict set gives the
open endpoint freedom required in the area argument.  The switch
roots, frames, and endpoint Weyl elements take only finitely many
values; pass to a subsequence on which these discrete data are fixed.
The ordered switch times lie in a compact simplex.  Uniform convergence
of the shapes then gives convergence of the positions and switch
coefficients.  The roots remain the same independent basis, and the
signs become weak.  Central coordinates have been carried unchanged
by addition throughout.  This proves the stated limiting schedule.
\end{proof}

We record the concrete consequence of the cover property needed in
the integrality argument.  For type $D$, use signed labels $a$ with
$y_a=\pm y_{|a|}$ and likewise for the moving values $w\mu$.

\begin{corollary}[No strict third-axis obstruction]\label{real:blocker}
At a switch of the limiting schedule in
Proposition~\ref{prop:basic-schedule}, suppose the two exchanged
signed labels have strict opposite position and moving-value orders.
No signed label on a third physical axis can lie strictly between
them in both orders.  The moving values are read immediately before
the switch, including its specified position in a simultaneous cluster.
\end{corollary}

\begin{proof}
Such strict inequalities would persist in the regular approximating
schedules.  In their local sorted frame the third label would give
the two positive roots in \eqref{real:inversion-pair}, both inverted
by the incoming direction.  The switch would then have Bruhat length
drop greater than one, contrary to its cover property.
\end{proof}
